The inside part of the planetarium round consists of 3 questions and takes
30 minutes.

\begin{description}
\item[Question 1.] The sky projected in the dome corresponds to Suceava
  (longitude $26^\circ 15'$), at 18:00 UT, on a certain day of a certain
  month. Two arcs of circle will now be projected. The arcs are segmented.
  Each segment represents an interval of some degrees. This number is not
  the same for each arc.
  \begin{description}
  \item[(a)] Identify each arc by circling the correct name (Equator /
    Meridian / Ecliptic) and give the angular size of each segment (in
    degrees).
    [Answer table: first arc, second arc --- name and segment size]
  \item[(b)] Estimate the local sidereal time of the sky you see in the
    dome.
  \item[(c)] Determine the month to which the projected sky would correspond
    at the given time. Fill in the box the number of the month (1 to 12).
  \end{description}
\item[Question 2.] The assistant will use a small red arrow pointer to point
  at objects in the sky; each object will be pointed at for 2 minutes.
  Locations of three Messier objects will be pointed at one by one. For each
  Messier object pointed at, fill in the boxes the Messier catalogue number
  and the number which indicates its type, as follows: 1 for galaxy, 2 for
  nebula, 3 for open cluster, 4 for globular cluster. Also, for each object,
  fill in the appropriate box the IAU abbreviation of the constellation
  where the object is located, using the provided table of constellations.
  [Answer table: 1st, 2nd, 3rd Messier object --- number, type,
  constellation]
\item[Question 3.] Three stars will be pointed at successively; each star
  will be pointed at for 2 minutes. Fill in the appropriate box the name of
  the star (or Bayer designation) and the number which indicates its type
  (1 for single, 2 for double). Also, for each star, fill in the appropriate
  box the IAU abbreviation of the constellation where the star is located,
  using the provided table of constellations.
  [Answer table: 1st, 2nd, 3rd star --- name, type, constellation]
\end{description}
